
\subsubsection{Key Findings}

\textbf{Finding 1:} The popularity-gap correlation is substantial (21.21 pp). Models trained on CIFAR-10 exhibit generalization gaps that are categorically different from models trained on SVHN---the former degrades by nearly 19 percentage points while the latter improves by over 2 percentage points on held-out data.

\textbf{Finding 2:} Research investment is highly concentrated. High-use datasets receive 7.56 times more optimization-focused papers, consistent with the hypothesis that popular benchmarks attract disproportionate research attention.

\textbf{Finding 3:} Texture bias does not explain the effect. Contrary to expectations from prior work on ImageNet, ResNet-18 shows \textit{lower} texture bias than VGG-11 at CIFAR scale. This suggests that architectural innovations (skip connections, batch normalization) improve rather than degrade shape-based generalization.

\subsubsection{Mechanistic Interpretation}

The verified causal chain is:

\begin{verbatim}
[High Popularity] → [More Optimization Research (7.56x)] → [???] → [Larger Generalization Gap]
\end{verbatim}

Step 1 (research investment differential) is verified. Step 2 (texture bias mechanism) is refuted. The pathway from intensive optimization to generalization failure remains unknown but is not mediated by texture bias at this scale.

Possible alternative mechanisms include:
\begin{itemize}
\item Spurious correlation exploitation (background features, shortcuts)
\item Test set leakage through extensive community usage
\item Architecture-specific inductive biases matching dataset statistics
\end{itemize}

\subsubsection{Limitations}

\begin{enumerate}
\item \textbf{Two dataset pairs:} The study cannot claim generality across all domains. CIFAR-10 (natural images) and SVHN (digits) represent different visual domains; the observed gap difference may partially reflect domain-specific characteristics rather than popularity alone.
\end{enumerate}

\begin{enumerate}
\item \textbf{Single-run proof-of-concept:} Statistical significance for H-E1 requires multi-run replication. The effect size (21.21 pp) is large enough that the direction is unambiguous, but formal significance testing is not available.
\end{enumerate}

\begin{enumerate}
\item \textbf{Incomplete mechanism:} Texture bias is ruled out; alternative mechanisms (spurious correlations, test set leakage) remain untested.
\end{enumerate}

\begin{enumerate}
\item \textbf{Bibliometric proxy:} arXiv paper counts serve as a proxy for optimization investment. Actual hyperparameter search compute is not directly observable.
\end{enumerate}

\begin{enumerate}
\item \textbf{Architecture-era confound:} The H-M2 comparison between VGG-11 and ResNet-18 cannot isolate ''optimization intensity'' from ''architectural innovation.'' The lower texture bias in ResNet may reflect skip connection design rather than optimization history.
\end{enumerate}

\subsubsection{Broader Impact}

These findings suggest that benchmark selection itself may be a confounding factor in model evaluation. Practitioners should consider validating on less-popular, less-optimized datasets before deployment. The research community should consider whether progress measured on popular benchmarks reflects true generalization capability.

